\documentclass[twocolumn]{article}
\usepackage[a4paper, margin=1cm]{geometry}

\usepackage[small]{titlesec}
\usepackage{url}
\usepackage{graphicx}
\usepackage{amsmath}
\usepackage{physics}
\usepackage{mathtools}
\usepackage{bbold}
\usepackage{float}
\usepackage{array}

\usepackage{lmodern}
\usepackage[T1]{fontenc}
\usepackage[fontsize=11pt]{fontsize}

\DeclareMathOperator{\erfi}{erfi}
\newcommand{\R}{\ensuremath{\mathbb{R}}}
\newcommand{\C}{\ensuremath{\mathbb{C}}}

\title{Complex Systems Excercise 5}
\author{Idan Stark}
\date{\today}

\begin{document}

\maketitle

\section{Dimensional Analysis}
\subsection{The pi theorem}
Every physical law that relates $n$ physical quantities containing $k$ independent dimensions
can be expressed as a relation between $p = n - k$ dimensionless variables, which are
often noted $\pi$.
\subsection{What to do?}
Steps (or comments? Or general things to do?):
\begin{itemize}
    \item Write the dimension of each quantity.
    \item Find the smallest set of independent dimensions (default to SI, but attempt complex dimensions as well).
    \item Write down the dimensionless variables.
    \item Write one dimensionless variable as a function of the others. Don't forget to add a function for each dimensionless quantity.
    \item Get a self similar form: an equation for a dimensionless variable.
    \item Solve it and go back to initial variables.
\end{itemize}
\section{Continuum Mechanics}
\subsection{Density, Velocity, and Compressability}
Continuity equation
\begin{equation}
    \partial_t\rho + \nabla\cdot(\rho v) = 0
\end{equation}
\textbf{Incompressability}
\newline
Either $\rho = \rho_0$ is a constant in space and time, or (the more common definition):
\begin{equation}
    \nabla\cdot v = 0
\end{equation}
For a incompressible fluid (by the second definition), the Continuity equation becomes:
\begin{equation}
    \partial_t\rho + v\cdot\nabla\cdot\rho = 0
\end{equation}
\textbf{Lagrangian Coordinates (or material coordinates, or convected coordinates)}
\newline
These are coordinates that move with the fluid. The material derivative is defined as:
\begin{equation}
    D_t = \partial_t + v\cdot\nabla
\end{equation}
Given a regular (Eulerian) coordinate $x$, then $X$ is a Lagrangian Coordinate if:
\begin{equation}
    D_t f = \frac{\partial f}{\partial t}\Bigg|_X
\end{equation}
From the Continuity equation:
\begin{equation}
    D_t \rho = -\rho \nabla \cdot v
\end{equation}
So for a incompressible fluid the material derivative of $\rho$ is zero.
\subsection{Conservation of Momentum}
\begin{equation}
    D_t \boldsymbol{v} = \boldsymbol{F} + \frac{1}{\rho}\boldsymbol{\nabla}\sigma
\end{equation}
Where $\boldsymbol{F}$ is the force per unit mass and $\sigma_{ij}$ is $i$ component the force per unit area with normal in the $j$th direction.
\newline
\textbf{Flow by inertia}
When $F = 0$ and $\sigma = 0$:
\begin{equation}
    D_t \boldsymbol{v} = 0
\end{equation}
Solved (in 1D) by $v(x, t)$ where:
\begin{equation}
    v = f(\xi)
    \qquad
    \xi = x - vt
    \qquad
    f(x) = v(x, 0)
\end{equation}
\begin{equation}
    \partial_x v = \frac{f'}{1+tf'}
    \qquad
    \partial_t v = -\frac{vf'}{1+tf'}
    \qquad
    D_t \xi = 0
\end{equation}
For $f' < 0$ there exists a time
\begin{equation}
    t^* = \frac{1}{f'} = \frac{1}{v^*}
\end{equation}
Where the velocity gradient $\partial_x v$ diverges.
\subsection{Ideal Incompressable Flow}
For a fluid with isotropic pressure $\sigma_{ij} = -\delta_{ij}p$:
\begin{equation}
    D_t \boldsymbol{v} = \boldsymbol{F} - \frac{1}{\rho}\boldsymbol{\nabla}p
\end{equation}
Assuming Incompressability $\nabla\cdot v = 0$ and conservative forces $\nabla\cross F = 0$, and taking the curl gives:
\begin{equation}
    \partial_t \omega + v\cdot\nabla \omega + \omega\cdot\nabla v = 0
\end{equation}
Where $\omega = \nabla \cross v$ is the \textbf{vorticity}. Lagrange's theorem states that if $\omega(t = 0) = 0$ then $\omega = 0$ in all times.
\newline
\textbf{Potential flow}
If $\omega = 0$, then $v = \nabla \phi$ for some potential function $\phi$ and compressability is $\nabla^2 \phi = 0$.
\subsection{Bernoulli's Principle}
Assuming Incompressability and that forces are gravitaitonal $\textbf{F} = -g\hat{\textbf{z}}$:
\begin{equation}
    \frac{1}{2}\rho v^2 + P + \rho gz = const
\end{equation}
 If the $\omega = 0$ (potential flow), this is true everywhere. If $\omega \ne 0$, this is true \textbf{along streamlines} (constant Lagrangian coordinates).
\subsection{Viscous flow}
Symmetry of the stress tensor:
\begin{equation}
    \sigma_{ij} = \sigma_{ji}
\end{equation}
Gives the \textbf{Navier Stokes equations}:
\begin{equation}
    \rho D_t v = \rho \partial_t \boldsymbol{v} + \rho(\boldsymbol{v}\cdot\boldsymbol{\nabla})\boldsymbol{v} = \rho \boldsymbol{F} - \boldsymbol{\nabla}p + \eta \nabla^2 \boldsymbol{v} + (\mu + \eta)\boldsymbol{\nabla}(\nabla\cdot v)
\end{equation}
Assuming incompressability and naming kinematic viscosity $\nu = \mu/\rho$:
\begin{equation}
    \partial_t \boldsymbol{v} + (\boldsymbol{v}\cdot\boldsymbol{\nabla})\boldsymbol{v} = \boldsymbol{F} - \frac{1}{\rho}\boldsymbol{\nabla}p + \nu \nabla^2 \boldsymbol{v}
\end{equation}
\textbf{Poiseuille flow}
\newline
Steady flow $\partial_t v = 0$, of a viscous incompressible fluid $(\rho, \eta)$, in a round pipe with radius $R$, with no external forces, flowing only along the pipe $\partial_\phi v = \partial_r v = 0$, assuming constant pressure gradient in the flow direction $\partial_x p = -G$:
\begin{equation}
    \boldsymbol{\nabla}p = \eta \nabla^2 v_x = \frac{\eta}{r}\partial_r(r\partial_r v_x(r))
\end{equation}
Requiring no divergences at $r = 0$ and $v_x(R) = 0$:
\begin{equation}
    v_x = -\frac{\Delta p}{4\eta L}\left(R^2 - r^2\right)
\end{equation}
And the mass flux:
\begin{equation}
    Q = 2\pi \int_0^R v_x rdr = \frac{\pi R^4}{8\eta}\frac{-\Delta p}{L}
\end{equation}
\textbf{Reynold number}
\newline
Dimensionless number describing the reltions in NS equation.
\begin{equation}
    Re = \frac{lv}{\nu} = \frac{lv\rho}{\eta}
\end{equation}
Where $l, v$ are the typical length scale and velocity.
In Navier-Stokes:
\begin{equation}
    \partial_t \boldsymbol{v} + (\boldsymbol{v}\cdot\boldsymbol{\nabla})\boldsymbol{v} = \boldsymbol{F} - \frac{1}{\rho}\boldsymbol{\nabla}p + \frac{1}{Re} \nabla^2 \boldsymbol{v}
\end{equation}
\begin{itemize}
    \item $Re \gg 1$ - the system is dominated by inertia
    \item $Re \ll 1$ - inertia is negligable (Stokes flow)
\end{itemize}
\textbf{Stokes flow}
\newline
Neglecting the inertia terms:
\begin{equation}
    \boldsymbol{F} - \frac{1}{\rho}\boldsymbol{\nabla}p + \nu \nabla^2 \boldsymbol{v} = 0
\end{equation}
So, for an object with length scale $r_0$ (like a sphere), using dimensional analysis
\begin{equation}
    F = C \eta r_0 v
\end{equation}
For sphere, $C = 6\pi$.
\newline
\textbf{The Oseen Tensor}
\newline
The Green function for Stokes flow equation. For a force distribution $f(r)$:
\begin{equation}
    p(r) = \int g(r - r')f(r')d^3r'
\end{equation}
\begin{equation}
    v(r) = \int T(r - r')f(r')d^3r'
\end{equation}
We get:
\begin{equation}
    g_i = \frac{r_i - r_i'}{4\pi|\boldsymbol{r} - \boldsymbol{r}'|^3}
\end{equation}
\begin{equation}
    T_{ij} = \frac{1}{8\pi|\boldsymbol{r} - \boldsymbol{r}'|}\left(\delta_{ij} + \frac{(r_i - r_i')(r_j - r_j')}{|\boldsymbol{r} - \boldsymbol{r}'|^2}\right)
\end{equation}
\subsection{Hydrodynamics Instabilities}
The instability of two fluid with relative tangential velocity along their surface x-y. Since in Navier Stokes for a compressible fluid, the pressure follows Laplace's equation $\nabla^2p = 0$, and there can be pressure waves:
\begin{equation}
    p = p_0 + \delta p \qquad \delta_p = Ae^{i(kx - \omega t) - kz}
\end{equation}
Allowing for velocity waves as well ($v$ is the relative velocity of $\rho_1$ with respect to $\rho_2$):
\begin{equation}
    v = v_0 + \delta v \qquad \delta_v = Be^{i(kx - \omega t) - kz}
\end{equation}
solving Navier-Stokes, and requiring that the fluids interface $z = \zeta(x, t)$ has the same velocity gives:
\begin{equation}
    \omega = kv_0\frac{\alpha \pm i\sqrt{\alpha}}{1 + \alpha} \qquad \alpha = \frac{\rho_1}{\rho_2}
\end{equation}
Which has an imaginary part, marking instability.
\newline
\textbf{Rayleigh Plateau instability}
\newline
The instability of a surface given a wave in it. In a cylinder of fluid with radius (with a volume preserving perturbation):
\begin{equation}
    r = r_0 + \epsilon \cos(kz) - \frac{\epsilon^2}{4\pi}
\end{equation}
Will have an area:
\begin{equation}
    A = \frac{4\pi^2}{k}\left(r_0 + \frac{\epsilon^2}{4\pi r_0}\left(r_0^1k^2 - 1\right)\right)
\end{equation}
Given the surface tension energy $E = \gamma A$, the energy is negative for $k < 1/r_0$, meaning that the perturbation will grow.
\subsection{Surfaces}
\subsection{Surface Tension}
\begin{equation}
    E = \gamma A = \int \sqrt{|a|}dx^1dx^2
\end{equation}
For an open surface, looking for minimum to the surface:
\begin{equation}
    H = 0
    \qquad
    -\frac{1}{2\nabla^2}\eta + \eta K > 0 \ \forall \ \eta
\end{equation}
For a colsed surface there can be a finite pressure difference $\delta p$, using \textbf{Young-Laplace law}.
\begin{equation}
    \Delta p = -2\gamma H
\end{equation}
\section{Elasticity}
The momentum equation (similar to NS):
\begin{equation}
    \partial_t v_i + v_j \partial_j v_i = f_i + \frac{1}{\rho}\partial_j \sigma_{ij} 
\end{equation}
For a static system $v_i = 0$:
\begin{equation}
    \partial_j \sigma_{ij} = 0
\end{equation}
The linearized strain tensor for some displacement $\boldsymbol{u} = \boldsymbol{r} - \boldsymbol{x}$, where at $\boldsymbol{x}$ the system is at rest:
\begin{equation}
    \epsilon_{ij} = \frac{1}{2}\left(\partial_i u_j + partial_j u_i\right)
\end{equation}
Their relation is:
\begin{equation}
    \sigma_{ij} = \lambda\delta_{ij}\Tr(\epsilon) + 2\mu \epsilon_{ij}
\end{equation}
Where $\lambda, \mu$ are the \textbf{Lame coefficients}.
\newline
\textbf{Young's modulus and Poisson ratio}
\newline
\begin{equation}
    \lambda = \frac{Y\nu}{(1 + \nu)(1 - 2\nu)}
    \qquad
    \mu = \frac{Y}{2(1 + \nu)}
\end{equation}
Young's modulue has units of pressure, and Poisson ratio is unitless, between $-1$ and $\frac{1}{2}$.
\subsection{Compression waves in solids}
Length wave $\epsilon_{xx} = \partial_x u_x$, gives
\begin{equation}
    \sigma_{ij} = (\lambda\delta_{ij} + 2\mu \delta_{ix}\delta_{jx})\partial_x u_x
\end{equation}
\begin{equation}
    \partial_t^2 u_x = C_p^2\partial^2 u_x
    \qquad
    C_p^2 = \frac{Y}{\rho}\frac{(1 - \nu)}{(1 + \nu)(1 - 2\nu)}
\end{equation}
Trasverse wave $\epsilon_{xy} = \partial_x u_y$, gives
\begin{equation}
    \sigma_{xy} = \sigma_{yx} = \mu \partial_x u_y
\end{equation}
\begin{equation}
    \partial_t^2 u_y = C_s^2\partial^2 u_y
    \qquad
    C_s^2 = \frac{Y}{\rho}\frac{1}{2(1 + \nu)}
\end{equation}
\subsection{Large displacements}
Strain to second order:
\begin{equation}
    \epsilon_{ij} = \frac{1}{2}\left(\partial_i u_j + \partial_j u_i + \partial_i u_j \partial_j u_i\right)
\end{equation}
Cauchy-Green strain:
\begin{equation}
    \epsilon = \frac{1}{2}\left(g - \bar{g}\right)
\end{equation}
Where $g$ is the metric and $\bar{g}$ is the original metric.
\subsection{This structure elasticity}
The 3D elastic energy functional:
\begin{equation}
    E = \int \frac{1}{2}\sigma^{ij}\epsilon_{ij}d^3x
\end{equation}
For thin structure $t, W \ll L$:
\begin{equation}
    E = YW\int \frac{1}{2}(\epsilon^1_1)^2 d^2x
\end{equation}
Thin configuration along a curve:
\begin{equation}
    \boldsymbol{r}(u, v) = \boldsymbol{\gamma}(u) + v\hat{\boldsymbol{n}}(u)
    \qquad
    \partial_u \boldsymbol{\gamma} = \frac{\partial s}{\partial u}\hat{\boldsymbol{t}}
\end{equation}
Where $\hat{\boldsymbol{t}}$ is the tangent and $\hat{\boldsymbol{n}}$ is the normal:
\begin{equation}
    \partial_u \hat{\boldsymbol{t}} = \frac{\partial s}{\partial u} \kappa \hat{\boldsymbol{n}}
    \qquad
    \partial_u \hat{\boldsymbol{n}} = -\frac{\partial s}{\partial u} \kappa \hat{\boldsymbol{t}}
\end{equation}
Assuming a small displacement $\frac{\partial s}{\partial u} = 1 + \Delta$, the metric is:
\begin{equation}
    a = \begin{pmatrix}
        (1 - \nu k)^2(1 + \Delta)^2 & \\ & 1
    \end{pmatrix}
\end{equation}
\textbf{The Kirchoff-Love assumptions}:
\newline
From the metric we get $\epsilon_{22} = 0$, and before we assumes $\sigma_{22} = 0$. These contradict. This finally gives:
\begin{equation}
    E = \frac{YWt}{2}\int_0^L \Delta^2 du + \frac{YWt^3}{24}\int_0^L \kappa^2 du \equiv E_{strech} + E_{bending}
\end{equation}
\section{Random Walks}
\textbf{In 1D}
\begin{equation}
    x_{n+1} = x_n \pm 1
\end{equation}
After $N$ steps, $R$ steps to the right and $N_R$ steps to the left:
\begin{equation}
    x_N = 2R - N
\end{equation}
\begin{equation}
    \langle x_N \rangle = 0 \qquad
    \langle x_N^2 \rangle = N
\end{equation}
\begin{equation}
    p(x) = \begin{cases}
        0 & x \ \text{odd} \\
        \frac{1}{2^N}\binom{N}{R} & x \ \text{even}
    \end{cases}
\end{equation}
Using Stirling's formula:
\begin{equation}
    p_{even}(x) = \frac{2}{\sqrt{2\pi N}}e^{-\frac{x^2}{2N}}
\end{equation}
\textbf{In higher dimensions}
\newline
Given $p(\vec{r}, t)$, we can write (time is still discrete, just with time steps $\tau$):
\begin{equation}
    p(\vec{r}, t + \tau) = \int p(\vec{r} - \vec{Delta}, t) g(\vec{Delta}) d^3\vec{Delta}
\end{equation}
Assuming isotropic distirbution and definint $\partial_t p = \frac{1}{\tau}(p(t + \tau) - p(t))$:
\begin{equation}
    \partial_t p = D \nabla^2 p
    \qquad
    D = \frac{1}{6d\tau}\langle \Delta^2 \rangle
\end{equation}
This is the diffusion equation, with Gaussian Solutions:
\begin{equation}
    p(\vec{r}, t) = \frac{1}{\left(4\pi D t\right)^{3/2}} e^{-\frac{r^2}{4Dt}}
\end{equation}
\begin{equation}
    \langle r^2 \rangle = 6Dt    
\end{equation}
\section{Markov Chains}
Given the probability matrix $P_{ij}$: the probability to go from state $i$ to state $j$, then the probability to be in state $i$ at $n+1$ is:
\begin{equation}
    p_i^{n+1} = \sum_j P_{ji}p_j
\end{equation}
So:
\begin{equation}
    \vec{p}_n = P^n \vec{p}_0
\end{equation}
\section{Langevin and Fokker-Planck Equations}
\textbf{Langevin equations} are the equations of motion with noise, and \textbf{Fokker-Planck equations} are the equations for the probability distribution.
\newline
For the simple equation
\begin{equation}
    \dot{x} = \xi(t)
\end{equation}
With $\langle \xi(t)\xi(t') \rangle = \frac{D}{2}\delta(t - t')$, then:
\begin{equation}
    \langle x^2 \rangle = 2Dt
\end{equation}
For the general Langevin equation:
\begin{equation}
    m\frac{dv}{dt} = -\gamma v + \xi(t) - V'(x)
\end{equation}
In the overdamped limit (removing the LHS and dividing by $\gamma$):
\begin{equation}
    \frac{dx}{dt} = \mu f(t) + \xi(t)
\end{equation}
We get two currents - force current and diffusion current:
\begin{equation}
    \vec{j}_f = -\mu p \vec{\nabla}V(\vec{r})
    \qquad
    \vec{j}_d = -D \vec{\nabla}p
\end{equation}
So from the Continuity equation:
\begin{equation}
    \partial_t p = -\vec{\nabla}\cdot\left[\vec{j}_f + \vec{j}_d\right]
\end{equation}
\section{Generalized CLT}
Characteristic function:
\begin{equation}
    \varphi(\omega) = \int_{-\infty}^\infty f(t)e^{i\omega t}dt
\end{equation}
For sum of variables $X = \sum_i x_i$:
\begin{equation}
    \varphi_X(\omega) = \prod_i \varphi_i(\omega)
\end{equation}
So:
\begin{equation}
    p_n = \mathcal{F}^{-1}\left[\phi_n\right] = \mathcal{F}^{-1}\left[\left(\mathcal{F}\left[p\right]\right)^n\right]
\end{equation}
For a distribution whose tail goes as $p(x > x^*) \propto x^{-(1 + \mu)}$:
\begin{itemize}
    \item $1 < \mu < 2$: the mean is defined but the variance isn't.
    \item $0 < \mu < 1$: the mean is not defined.
\end{itemize}
\textbf{Levi stable distribution} is a distribution $f(x)$ such that the sum of two random variables from the distribution has the same distribution. The most general such distribution has the characteristic function:
\begin{equation}
    \varphi(\omega) = \exp[A\left|\omega\right|^\mu]
\end{equation}
\section{Extreme Value Theorem}
The Cumulative Probability Function (CPF) of a maximum value scales multiplicativly:
\begin{equation}
    G(x) = p(X_n < x) = \prod_i p(x_i < x)
\end{equation}
For a distribution normalized using
\begin{equation}
    \xi_n = \frac{x_n - b_n}{a_n}
\end{equation}
Where
\begin{equation}
    \alpha = \lim_{n\rightarrow\infty} \frac{n}{a_n}\frac{\partial a_n}{\partial n}
    \qquad
    \beta = \lim_{n\rightarrow\infty} \frac{n}{a_n}\frac{\partial b_n}{\partial n}
\end{equation}
Then:
\begin{itemize}
    \item $\alpha > 0$ \textbf{Fréchet Distribution} has $G(x) = e^{-ax^{-1/\alpha}}$
    \item $\alpha = 0$ \textbf{Gumbel Distribution} has $G(x) = e^{-e^{-ax + b}}$
    \item $\alpha < 0$ \textbf{Weibull Distribution} has $G(x) = e^{\alpha(\beta/|\alpha| - x)^{1/|\alpha|}}$
\end{itemize}
\section{Random Matrix Theory}
There are three meaningful ensembles:
\begin{itemize}
    \item \textbf{Gaussian Orthogonal Ensemble}: matrices are real and symmetric, with probability $p(A) \propto e^{-\frac{N}{4}\tr A^2}$.
    \item \textbf{Gaussian Unitary Ensemble}: matrices are complex and hermitian, with probability $p(A) \propto e^{-\frac{N}{2}\tr A^2}$.
    \item \textbf{Gaussian Symplectic Ensemble}: matrices are complex and harmitian, with probability $p(a) \propto e^{-n\tr A^2}$
\end{itemize}
\textbf{Wigner's Surmise} says that the result of $N=2$ matrix are a good approximation to the $n \gg 1$ case.
\section{Dynamical Systems}
\subsection{Discrete maps}
\begin{equation}
    \vec{x}_{n+1} = f\left(\vec{x}_n\right)
\end{equation}
Fixed point:
\begin{equation}
    \vec{x}^* = f\left(\vec{x}^*\right)
\end{equation}
GRaphically these are the intersection of the map and the $y = x$ line.
\newline
Periodic orbits:
\begin{equation}
    \vec{x}^* = f^n\left(\vec{x}^*\right)
\end{equation}
Stability of a p-orbit is determined by its \textbf{Stability coefficient}:
\begin{equation}
    \lambda_p = \prod_i^p f'(x_i)
    \qquad
    \begin{matrix}
        \left|\lambda_p\right| > 1 & \text{unstable} \\
        \left|\lambda_p\right| < 1 & \text{stable} \\
        \left|\lambda_p\right| = 1 & \text{superstable}
    \end{matrix}
\end{equation}
Bifurcations:
\begin{itemize}
    \item Tanget bifurcation: two points, stable and unstable, appear/annahilate when the map $f(x_n)$ becomes tangent to $y = x$. Inlucdes a slowing of the dynamical system before the the creation (Intermittency).
    \item Period Doubling bifurcation: a stable fixed point becomes unstable and creates two stable 2-orbits.
    \newline
    This happens when the position of the fixed point is not changing, but $f'(x^*)$ crosses $1$ or $-1$. When crossing $1$, this is very similar to 1D continuous system Pitchfork bifurcation. When crossing $-1$, then $f^{2}{'}(x^*)$ crosses $1$, and gets two stable fixed point of $f^2$ are created. this essentially doubles the orbit of the 1-orbit to 2-orbit, and if $f$ is actually $f^p$, then it doubles the p-orbit to 2p-orbit.
    \item Inverse Period Doubling bifurcation: the same as period doubling, but a single unstable fixed-point (p-orbit) becomes two unstable 2-orbits (2p-orbits) and a stable fixed point.
\end{itemize}
\textbf{Sarkovskii's theorem}
\newline
Given Sarkovskii's order:
\begin{equation}
    3,5,7,\dots,2\cdot3,2\cdot5,\dots,2^2\cdot3,2^2\cdot5,\dots,2^3,2^2,2,1
\end{equation}
Then if a map has p-periodic orbit, it has periodic orbits of all order above it.
\subsection{Continuous time}
\begin{equation}
    \dot{\vec{x}} = f(t, \vec{x})
\end{equation}
Fixed point:
\begin{equation}
    f(\vec{x}) = 0
\end{equation}
\subsubsection{1D systems}
Categorization of fixed points:
\begin{itemize}
    \item $f'(x^*) < 0 \rightarrow$ stable
    \item $f'(x^*) > 0 \rightarrow$ unstable
    \item $f'(x^*) = 0 \rightarrow$ look at higher derivatives. if all derivatives are zero, the point is half-stable.
\end{itemize}
The characteristic time scale:
\begin{equation}
    \tau = \left|f'(x^*)\right|^{-1}
\end{equation}
\textbf{Bifurcations}
\begin{itemize}
    \item Saddle-Node bifurcation ($r + x^2$): when two fixed points join to become a half stable fixed point, then vanish.
    \item Transcritical bifurcation ($rx + x^2$): when a node changes its stability.
    \item Pitchfork bifurcation - in a symmetric system, a fixed point is split into three - one unstable and two stable fixed points.
    \newline
    These can be \textbf{supercritical} ($rx - x^3$), where the $x^3$ term causes the decay to be algebric with time (critical slowdown) or \textbf{subcritical} ($rx + x^3 - x^5$), where there are additional two stable fixed points, created by two simultanious saddle-point bifurcations, and $x = 0$ is only locally stable.
\end{itemize}
\subsubsection{2D systems}
\textbf{Stability}
\begin{itemize}
    \item Attracting fixed point: points in the neighborhood go to the fixed point at $t\rightarrow\infty$.
    \item Lyapunov stablity - If a point in the neighborhood remains in the neighborhood for all times $t$.
    \item A point that has stability in one direction and instability in the other is called a saddlepoint. It has a stable manifold (all points that flow to the fixed point at $t\rightarrow\infty$) and unstable manifold (the same at $t\rightarrow-\infty$). 
\end{itemize}
\textbf{Linear system}
\begin{equation}
    \dot{\vec{x}} = A\vec{x}
\end{equation}
So $\vec{x} = (0,0)$ is always a fixed point.
Given that $\lambda_{1,2}$ are the eigenvalues of $A$ and $\vec{v}_{1,2}$ are the eigenvectors, then:
\begin{equation}
    \lambda_{1,2} = \frac{\tau \pm \sqrt{\tau^2 -4\Delta}}{2}
\end{equation}
Where $\tau$ is the trace and $\Delta$ is the determinant, and the general solution starting from $\vec{x}(0) = c_1\vec{v}_1 + c_2\vec{v}_2$ is:
\begin{equation}
    \vec{x}(t) = c_1 \vec{v}_1 e^{\lambda_1 t} + c_2 \vec{v}_2 e^{\lambda_2 t}
\end{equation}
Then the point can be written as:
\begin{itemize}
    \item $\Delta < 0$ - saddlepoint: one eigenvalue is negative and one is positive.
    \item $\Delta > 0$ and $\tau^2 > 4\Delta$ - node: then the eigenvalues have the same sign. If $\tau > 0$ then they are both positive, and the node is unstable. If $\tau < 0$ then the point is stable.
    \item $\Delta > 0$ and $\tau^2 < 4\Delta$ - spirals: the eigenvalues are complex, and based on the sign of $\tau$ can have either stability.
    \item $\Delta > 0$ and $\tau = 0$: centers - the eigenvalues are purely imaginary.
    \item $\tau^2 = 4\Delta$: the eigenvalues are the same, and we get a star (with stability based on $\tau$) if the multiplicity is algebric and a degnerate node, with a single direction, if the multiplicity is geometric.
\end{itemize}
\textbf{Nonlinear system}
Given:
\begin{equation}
    \dot{x} = f(x, y) \qquad
    \dot{y} = g(x, y)
\end{equation}
We linearize around a fixed point:
\begin{equation}
    \dot{x} = Jx
    \qquad
    J = \begin{pmatrix}
        \partial_x f & \partial_y f \\ \partial_x g & \partial_y g
    \end{pmatrix}
\end{equation}
If we get a marginal case - any line in the $\tau, \Delta$ graph, be cautious!
\newline
A \textbf{hyperbolic} point is a point where $\Re(\lambda) \ne 0$ for all eigenvalues. \textbf{The Hartman-Grobman theorem} states that locally, the phase near hyperbolic points are “topologically equivalent” to the linearized version.
\newline
Centers are nevertheless stable in \textbf{conservative systems} (systems with a conserved quantity, i.e. energy), and systems with \textbf{time reversal symmetry}.
\newline
\textbf{Index Theory}
\newline
For a system $\dot{\vec{x}} = \vec{f}(\vec{x})$ we define:
\begin{equation}
    \phi(x) = \atan(\frac{\dot{y}}{\dot{x}})
\end{equation}
Going along a closed loop $C$ in $\R^2$, the loop index is defined as the number of cycles $\phi$ does:
\begin{equation}
    I_C = \frac{1}{2\pi} \oint \dot{\phi}(x(\lambda))d\lambda = 
    \frac{1}{2\pi} \oint \left(\frac{\ddot{x}}{\dot{x}} - \frac{\ddot{y}}{\dot{y}}\right)d\lambda
\end{equation}
\begin{itemize}
    \item The loop always goes counterclockwise. The index doesn't depend on the direction of time.
    \item The contour can be smoothly deformed, without going through any fixed points, and not change its index.
    \item If $C$ is a trajectory, then $I_C = 1$.
    \item The index is $0$ if there are no fixed points inside the loop.
    \item This means that index of an isolated fixed point is the index of any loop that contains only itself.
    \item The index of a fixed point is independent of its stability.
    \item The index of nodes, degenerate nodes, spirals and centers is 1, and of saddles is $-1$.
    \item If a closed curve $C$ encloses multiple fixed points, its index the is the sum of their indices:
    \item Any closed orbit in state-space has an index IC = +1, and therefore it must enclose fixed points whose indices sum up to +1.
\end{itemize}
\textbf{Limit Cycles}
\newline
An isolated closed trajectory.
\newline
There are no limit cycles:
\begin{itemize}
    \item In potential systems $\dot{\vec{x}} = -\vec{\nabla}V$.
    \item If there is a Liapunov Function: a function $V(x)$ such that at the fixed point $V(x^*) = 0$ and everywhere else $V(x) > 0$ and $\dot{V}(x) < 0$.
    \item Dulac's Criterion: In an area where $\nabla \cdot (g\vec{f})$ doesn't change sign, for every $g(x)$.
\end{itemize}
\textbf{Poincare-Bendixson Theorem}
Consider a close bounded subset of the plane $R$, and that the vector field $\vec{f}(\vec{x})$ is smooth enough in some open subset of $R$. If $R$ does not have a fixed point, and if at least one trajectory never leaves $R$,
then $R$ contains a closed orbit.
\newline
To use this, we find a \textbf{attracting zone}: we plot the nullcline and find an area where the normal derivative to the boundary is always inward.
\newline
\textbf{Bifurcations}
Has all the 1D bifurcations, and also \textbf{Hopf bifurcations}.
\begin{itemize}
    \item When the eigenvalues don't have an imaginary part, it's the same as in 1D.
    \item If they have an imaginary part, and they change sign of the real value, it's a \textbf{supercritical Hopf bifurcation} - moving from stable spiral to an unstable spiral, creating a limit cycle.
    \item The \textbf{subcritical Hopf bifurcation} similar to 1D subcritical pitchfork, with term $x^3$ unstabalizes and an additional $x^5$ stablizing.
    \item It is impossible to tell the type of Hopf bifurcation from linearization.
\end{itemize}
% quazi
\section{Math}
\subsection{Tensors}
Symmetric rank rwo tensors in 3D, linear to a generic tensor $A_{ij}$:
\begin{equation}
    \sigma_{ij} = \alpha \delta_{ij} + \beta Tr(A)\delta_{ij} + \gamma (A_{ij} + A_{ji})
\end{equation}
Symmetric, uniform, isotropic, rank four tensor:
\begin{equation}
    A^{ijkl} = \lambda\delta^{ij}\delta^{kl} + \mu\left(\delta^{ik}\delta^{jl} + \delta^{il}\delta^{jk}\right)
\end{equation} 
\newline
\textbf{Isotropy}
When requiring isotropy in a vector relation:
\begin{itemize}
    \item Wite the relation in tensor form $C^{i} = T^{ijk}A_jB_k$
    \item Move to another system $\Lambda^i_j$:
    \begin{equation}
        C'^{i} = T^{jkl}\Lambda^i_j \Lambda^n_k \Lambda^m_l A_nB_m
    \end{equation}
    \item Reuiqre $T$ to be invariant
    \begin{equation}
        T^{inm} = T^{jkl}\Lambda^i_j \Lambda^n_k \Lambda^m_l
    \end{equation}
    \item Rotate by different tegrees to see relations between the components of $T$
\end{itemize}
In the case of the bilinear tensor relation $C^{i} = T^{ijk}A_jB_k$:
\begin{equation}
    T^{ijk} \propto \epsilon^{ijk}
\end{equation}
\subsection{Surfaces}
Given a 2D manifold in 3D space:
\begin{equation}
    \boldsymbol{r}(x^1, x^2): \mathcal{D}\rightarrow\mathcal{S}
\end{equation}
Where $\mathcal{D} \in \R^2$ and $\mathcal{S} \in \R^3$.
\newline
Tanget vectors:
\begin{equation}
    \partial_1 r, \partial_2 r
\end{equation}
Distances are measured using the metric (first fundamental form):
\begin{equation}
    ds^2 = g_{\alpha\beta}dx^\alpha dx^\beta
    \qquad
    g_{\alpha\beta} = \partial_\alpha r \partial_\beta r
\end{equation}
The normal vector:
\begin{equation}
    \hat{\boldsymbol{N}} = \frac{\partial r \cross \partial_2 r}{\sqrt{a}}
\end{equation}
The second fundamental form (the curvature form):
\begin{equation}
    b_{\alpha\beta} = \partial_\alpha \partial_\beta r \cdot \hat{\boldsymbol{N}}
\end{equation}
The shape tensor:
\begin{equation}
    s^\alpha_\beta = a^{\alpha\gamma}b_{\gamma\beta}
\end{equation}
The physical normal curvature of curve $(dx^1, dx^2)$
\begin{equation}
    k_n(dx^1, dx^2) = \frac{dx^\alpha b_{\alpha\beta}dx^\beta}{dx^\alpha a_{\alpha\beta}dx^\beta}
\end{equation}
For every given point, changing coordintes such that $a$ is the identity, then rotating such that $b$ is diagonal, gives:
\begin{equation}
    b = \begin{pmatrix}
        \kappa_1 & \\ & \kappa_2
    \end{pmatrix}
\end{equation}
Where $\kappa_1, \kappa_2$ are the \textbf{principal cuvatures}. They give the mean curvature and Gaussian curvature:
\begin{equation}
    H = \frac{\kappa_1 + \kappa_2}{2} = \frac{1}{2}Tr(s)
    \qquad
    K = \kappa_1\kappa_2 = |s| = \frac{|b|}{|a|}
\end{equation}
To get a minimal sruface:
\begin{equation}
    H = 0
    \qquad
    -\frac{1}{2\nabla^2}\eta + \eta K > 0 \ \forall \ \eta
\end{equation}
\subsection{Probabilities}
Stirling's formula:
\begin{equation}
    N! = \sqrt{2\pi N}\left(\frac{N}{e}\right)^N
\end{equation}
\subsection{Gaussian Integrals}
Simple Gaussian:
\begin{equation}
    \int dx  e^{-ax^2 + bx} = \sqrt{\frac{\pi}{a}}e^{\frac{b^2}{2a}}
\end{equation}
$N$-dimensional, with source $J$:
\begin{equation}
    \int d^n x  e^{-\frac{1}{2} x^T A x + J^T x} = \sqrt{\frac{(2\pi)^N}{\det A}} \exp\left( \frac{1}{2} J^T A^{-1} J \right)
\end{equation}
% TODO add Stokes first and second problem solutions at the end
\end{document}